\section{Scope, classification targets, and open boundaries}
\label{sec:publication-scope}

The global theorem is componentwise because the certified reversible core has five connected semantic doctrines.  The component tag is part of the invariant; omitting it would compare objects for which no reversible semantic bridge has been proved.

The Universal Lane Transport Theorem has a different scope.  It is an exact finite theorem inside a declared additive, idempotent-complete setting with complete split packet systems at a common resolution.  Its realization, kernel, multilinear, fold, and descent clauses apply only when their additional hypotheses are present.  Trivial address holonomy is required only to replace all local labels by one fixed global address set.  Under a coherent nontrivial finite address cocycle, complete matrices instead descend as compatible sections of a matrix local system, and monodromy orbits give the finest ordinary global projector coarsening.  The word ``universal'' refers to this formal lane calculus and its typed specializations.  It does not claim that every object possesses a finite lane presentation or that address count equals vector-space dimension, minimal rank, entropy, or information.

\input{generated/block14_scope_table.tex}

\subsection{Directional satellites}

Several mathematically useful constructions lie adjacent to the certified reversible core but are not equivalences.  The most important examples are:
\begin{itemize}
\item reduction from integral section lattices to finite residue systems, with Smith--Tor ghosts and missing saturated coordinates;
\item passage from an algebraic branch to a differential annihilator, which can forget the selected branch and its minimal polynomial;
\item functional-Mahler masks and Newton--Mahler interpolation, which use different variables and require additional lifting and interpolation data before comparison;
\item projective Galois descent, whose lift to an ordinary vector realization is obstructed by a relative Brauer class;
\item faithful additive or Hilbert-space realizations, which do not supply the exact tensor fiber functor required for Tannakian reconstruction.
\item incomplete lane selections and noninjective output folds, which retain explicit residual or collision kernels rather than becoming reversible transport.
\end{itemize}
These constructions remain in the directed transformation category with their exact loss or obstruction records.

\subsection{Open mathematical obligations}

The next mathematical obligations are narrower than the original foundational problem.
\begin{enumerate}
\item Extend the certified atlases across radices, coefficient rings, semirings, and mixed local-global completions without losing normalization or effectivity.
\item Develop modular and wild analogues of the finite semisimple Galois projector theory when Reynolds averaging fails.
\item Replace the marked model-family results in the holonomic, algebraic, and nonlinear strata by broader native minimality and collision invariants.
\item Construct nonneutral and integral tensor-descent variants with their gerbe or torsor data retained explicitly.
\item Prove finite-data recognition theorems only on declared subgrammars where identifiability and certification are available.
\item Extend the finite split and finite holonomy-twisted lane theorems to infinite, completed, nonsplit, noninvertible, or genuinely nonadditive resolutions only after specifying the required convergence, exactness, and coherence doctrines.
\end{enumerate}

\begin{boundarybox}
\textbf{No global promotion.}
Completion of the finite certified atlas, the finite profiled lane theorem, or its finite holonomy-twisted extension does not promote absent bridges, directional constructions, incompatible selectors, noninjective folds, or unresolved broad-stratum normal forms to equivalences.  The detailed obligation and claim ledgers are retained in \cref{app:block14-open-program,app:block14-claim-ledger}.
\end{boundarybox}
